\documentclass[11pt,a4paper]{article}
\usepackage[margin=2.3cm]{geometry}
\usepackage{booktabs}
\usepackage{graphicx}
\usepackage{amsmath}
\usepackage{microtype}
\usepackage{siunitx}
\sisetup{range-phrase={--},range-units=single}
\usepackage[colorlinks=true,linkcolor=black,urlcolor=blue!60!black]{hyperref}
\usepackage{xcolor}
\usepackage{caption}
\captionsetup{font=small,labelfont=bf}

\graphicspath{{../../results/structural/hanging/figures/}}
\newcommand{\billowmean}{15.7}
\newcommand{\kfemmean}{12.4}
\newcommand{\billowmeansix}{11.5}
\newcommand{\kfemmeansix}{7.8}
\newcommand{\billowtime}{14}
\newcommand{\billowspanone}{7.276}
\newcommand{\storedmean}{15.1}
\newcommand{\storedspanone}{7.212}
\newcommand{\kfemtime}{552}


\title{\textbf{Validating Billow against a hanging V3 kite}\\[2mm]
\large Measured deformed shape, ten load cases, and what the dataset can and
cannot decide}
\author{Structural module, AWETrim}
\date{\today}

\begin{document}
\maketitle

\begin{abstract}
\noindent
The TU~Delft V3 kite was suspended from its bridle and its deformed shape
measured under ten combinations of inflation pressure and applied load. This
report rebuilds that experiment in \textbf{Billow}, the standalone
minimum-energy structural model in \texttt{awetrim.structural}, and scores it
against both the measurements and the reference finite-element model
\texttt{kite\_fem}, on an element-for-element identical mesh.

Three things came out of it. First, the distributed reference model is
\emph{ill posed}: its anchor node belongs to no element, so the kite is a free
body under \SI{126}{N} of unbalanced weight, and the reference results are held
together by a numerical regularisation. Second, \texttt{kite\_fem}'s beams
\emph{converge} to about $0.23\times$ its own constitutive law, with negative
bending stiffness in several load cases --- so its closer agreement with the
measurement is not evidence of better physics. Third, and most limiting for the
original purpose, this dataset is governed almost entirely by inflatable-tube
bending and is nearly blind to the canopy model, so no canopy conclusion can be
drawn from it.

Billow reaches force balance to $\sim$\SI{1e-9}{N} in a mean of
\billowtime~seconds per case against the reference's \kfemtime~s. On the six
cases both codes can solve, the mean relative error is \billowmeansix\% for
Billow against \kfemmeansix\% for \texttt{kite\_fem}; over all ten it is
\billowmean\% against \kfemmean\%, the difference being four centre-load cases
that neither code reproduces --- the thesis says the same of its own model.
Neither reproduces the measured shape --- though the published data, six
independent numbers with no height or chordwise measure among them, is a weak
test of shape at all. The report closes with the sharpest
remaining clue: the measured pressure sensitivity is roughly five times what
the fitted tube law predicts, at pressures below the range over which that fit
was ever demonstrated.
\end{abstract}

\section{The experiment}

The kite hangs from its bridle and the deformed shape is recorded
stereoscopically, with markers on the inflatable structure. Twelve lengths are
extracted per load case: nine trailing-edge segments \texttt{La}--\texttt{Li}
between consecutive struts, the tip-to-tip span \texttt{b}, and two
tip-to-centre-leading-edge diagonals \texttt{LcsTL}, \texttt{LcsTR}.

The deformation is large. Under self-weight alone the span falls from
\SI{8.305}{m} as built to a measured \SI{4.867}{m} --- a 41\% closure --- and
every trailing-edge segment shortens by 26--43\%. The tip-load cases run the
same deformation backwards: spreading the tips with \SI{2}{kg} or \SI{5}{kg}
reopens the span to \SI{7.8}{}--\SI{8.3}{m}, close to the as-built shape.

\paragraph{Suspension.} The kite hangs from a \textbf{bar}, not a point.
Roeleveld's thesis (Appendix~B, Table~B.1) marks \texttt{A\_main},
\texttt{A\_I}, \texttt{M}, \texttt{PL} and \texttt{SL} as \emph{N/A} for the
hanging test and adds a \texttt{bar} of \SI{1650}{mm}; Section~6.3 explains
that the bridle was adapted to fit the hangar height. The same table records
the one documented shortening, \texttt{B\_r,5} from \SI{12168}{}
to~\SI{9425}{mm}.

\paragraph{Provenance caveat.} The upstream repository describes the
measurements only as ``experimental validation reference data'' and never uses
the word photogrammetry; the thesis identifies the method as stereoscopic
camera measurement. The published rows are also perfectly symmetric
left-to-right to four decimals, so they have been mirrored or
port/starboard-averaged.

\section{Rebuilding the model}

Billow is rebuilt from \texttt{hanging\_test\_initial.npz}, the fully assembled
reference model, rather than from the YAML --- there is no committed
YAML$\rightarrow$npz builder, so the \texttt{.npz} \emph{is} the input both
codes solve. Nodes, connectivity, rest lengths, stiffnesses and tube diameters
therefore transfer with no re-interpretation:

\begin{center}
\begin{tabular}{ll}
\toprule
shared, element for element & 276 nodes, 46 bridle cables, 14 pulleys,
                              98 inflatable tube beams \\
changed & canopy: 721 tension-only springs $\rightarrow$ 392 wrinkling
          membrane triangles \\
\bottomrule
\end{tabular}
\end{center}

Two checks make this trustworthy rather than assumed. The load vector
reproduces \texttt{kite\_fem}'s own \texttt{loading()} to \textbf{exactly zero}
for all five load combinations, and the length extractor reproduces its
published \texttt{model\_results.csv} from its saved states to
\textbf{machine zero}.

\paragraph{Canopy grid.} The reference canopy is a structured $29\times8$ grid
stored only as a flat spring list. It is recovered geometrically --- the
reference builder projects every interior section node onto the straight line
joining its leading- and trailing-edge node, so each chordwise section is
exactly collinear --- and then \emph{verified}: all 29 sections must come out
with 8 nodes, and all 721 canopy springs must be grid edges with no spring
joining two grid nodes being anything else. Both hold exactly.

\section{Defects found in the reference model}

\subsection{The anchor is connected to nothing}

Node~0 is flagged fixed but appears in no spring, no pulley and no beam; seven
further bridle nodes are orphaned with it (all of exactly zero mass). The kite
is therefore a free body under \SI{126}{N} of unbalanced weight, and its
potential is unbounded below. Billow reports that faithfully: the structure
translates hundreds of kilometres, and neither load stepping nor constitutive
continuation helps, because none of them supplies the missing constraint.

\texttt{kite\_fem} does not diverge only because its \texttt{I\_stiffness}
regularisation acts as a ground spring on every node. Its own published
solutions still drift \SIrange{0.29}{0.34}{m} bodily upward, with node~0 sitting
fixed and unconnected throughout.

\subsection{Bridle lengths disagree with the thesis}

Twenty-eight of the thirty-six declared bridle lines match Table~B.1 to the
millimetre, which is what makes the remaining eight look like transcription
errors rather than naming mismatches: \texttt{BRI} and \texttt{BRII} are an
exact swap; \texttt{abcd2} and \texttt{abcd3} carry a duplicated leading digit;
\texttt{ab3} equals \texttt{ab2}; \texttt{a5} and \texttt{br5} differ; and
\texttt{brmain ext} is present in the table but unwired.

The eighth is not a transcription error but a contradiction inside the thesis
itself. Table~B.1 marks \texttt{A\_I} as N/A for the hanging test, while
Figure~B.1 \emph{draws} \texttt{A\_I} running to the bar in the adapted
geometry. The two readings are structurally different --- cutting the line
removes a front load path --- so both were run:

\begin{table}[htbp]
\centering
\caption{Three readings of the same source. Cutting \texttt{A\_I} is the most
accurate reading on the cases the model can solve and the worst overall,
because removing that load path is what lets the four centre-load cases fall
into a folded local minimum. No reading is adopted here: the ambiguity is a
finding about the source data, not something to settle by taking whichever
scores best.}
\begin{tabular}{lrrrr}
\toprule
bridle reading & all ten & six solvable & four centre-load & folded \\
 & (\%) & (\%) & (\%) & cases \\
\midrule
bridle as stored & 15.1 & 11.0 & 21.3 & 0 \\
Table~B.1, $A_I$ cut & 17.3 & 8.9 & 30.0 & 4 \\
Table~B.1, $A_I$ kept & 15.7 & 11.5 & 21.9 & 0 \\
\texttt{kite\_fem} & 12.4 & 7.8 & 19.3 & 0 \\
\bottomrule
\end{tabular}
\end{table}

\subsection{The reference beams converge far softer than their own law}

Billow's integrated tube law agrees with \texttt{kite\_fem}'s secant $EI$ to
1--3\% across the whole curvature range, departing only below its
\num{0.03} normalised-deflection floor. The two codes therefore implement the
same constitutive law. They do not converge to the same stiffness:

\begin{center}
\begin{tabular}{rrrrr}
\toprule
load case & median $EI$ (\si{N.m^2}) & vs.\ own law & min $EI$ & beams flagged collapsed \\
\midrule
1 & 24.5 & $0.29\times$ & 2.7 & 78 \\
4 & 69.4 & $0.81\times$ & $-42.4$ & 49 \\
6 & 43.1 & $0.50\times$ & $-164$ & 76 \\
9 & 84.1 & $0.98\times$ & $-459$ & 40 \\
\bottomrule
\end{tabular}
\end{center}

Three of the four cases contain \emph{negative} bending stiffness, which the
secant form $EI = PL^3/\bigl(3(v - PL/kGA)\bigr)$ produces as soon as the
deflection estimate falls below the shear term. \texttt{kite\_fem} therefore
lands closer to the measurement than Billow because its beams are several times
softer than the law it implements, not because its physics is closer.

\section{Results}

\begin{table}[htbp]
\centering
\caption{All ten load cases. Errors are the mean over the twelve measured
lengths of $|L_\text{model}-L_\text{meas}|/L_\text{meas}$; the reference
\texttt{shapecorrelation} metric is not used, because as a dot-product ratio
dominated by the two quantities of order \SIrange{4}{8}{m} it reads
0.980--0.9997 even where the span is 48\% wrong. $\|r\|$ is the force residual
on the worst free node.}
\begin{tabular}{rccc rrr rr rr}
\toprule
 & \multicolumn{3}{c}{load case} & \multicolumn{3}{c}{span (m)} & \multicolumn{2}{c}{mean rel.\ error (\%)} & \multicolumn{2}{c}{Billow} \\
\cmidrule(lr){2-4}\cmidrule(lr){5-7}\cmidrule(lr){8-9}\cmidrule(lr){10-11}
LC & $p$ (bar) & tip (kg) & centre (kg) & meas. & Billow & \texttt{kite\_fem} & Billow & \texttt{kite\_fem} & $t$ (s) & $\|r\|$ (N) \\
\midrule
1 & 0.15 & 0 & 0 & 4.867 & 7.276 & 5.382 & 25.9 & 13.7 & 7 & 2.2e-07 \\
2$^{\dagger}$ & 0.15 & 0 & 9.7 & 5.042 & 7.554 & 7.462 & 27.6 & 25.0 & 15 & 9.1e-08 \\
3$^{\dagger}$ & 0.15 & 0 & 25.2 & 5.113 & 7.841 & 7.576 & 21.8 & 18.9 & 19 & 3.8e-07 \\
4 & 0.15 & 2 & 0 & 7.843 & 8.597 & 8.467 & 6.5 & 6.7 & 16 & 1.1e-07 \\
5 & 0.15 & 5 & 0 & 8.271 & 8.916 & 8.962 & 5.9 & 5.4 & 18 & 4.1e-07 \\
6 & 0.25 & 0 & 0 & 5.453 & 7.468 & 6.363 & 19.5 & 9.3 & 6 & 4.0e-09 \\
7$^{\dagger}$ & 0.25 & 0 & 9.7 & 5.575 & 7.690 & 7.691 & 19.7 & 18.9 & 15 & 3.7e-09 \\
8$^{\dagger}$ & 0.25 & 0 & 25.2 & 5.697 & 7.965 & 7.638 & 18.4 & 14.6 & 17 & 6.4e-09 \\
9 & 0.25 & 2 & 0 & 7.818 & 8.593 & 8.461 & 5.2 & 5.6 & 15 & 1.1e-07 \\
10 & 0.25 & 5 & 0 & 8.224 & 8.908 & 8.955 & 6.2 & 5.9 & 16 & 4.0e-07 \\
\midrule
\multicolumn{7}{r}{mean, all ten} & 15.7 & 12.4 & 14 & \\
\multicolumn{7}{r}{mean, excluding centre-load cases} & 11.5 & 7.8 & 13 & \\
\bottomrule
\end{tabular}
\end{table}

\begin{figure}[htbp]
\centering
\includegraphics[width=\textwidth]{summary.png}
\caption{Span and mean relative error across the ten cases. Both codes
overstate the span wherever they solve the physical branch; the gap is widest
where gravity alone sets the shape and narrowest where an imposed tip load
dominates. The shaded centre-load cases are reproduced by neither code, and
are where Billow's folded branch shows as a near-zero span.}
\end{figure}

\begin{table}[htbp]
\centering
\caption{Per-quantity error, averaged over the ten cases. The two codes fail
differently, which the signed-bias columns separate from the magnitude
columns: Billow's bias is one-signed --- it overstates essentially every
length, i.e.\ its wing stays too open --- whereas \texttt{kite\_fem} scatters,
understating the outboard segments and the tip-to-centre diagonals while
overstating the inboard ones. A single stiffness factor could plausibly
account for the first and could not account for the second.}
\begin{tabular}{lrrrrr}
\toprule
quantity & meas.\ range (m) & \multicolumn{2}{c}{$|$rel.\ error$|$ (\%)} & \multicolumn{2}{c}{signed bias (\%)} \\
\cmidrule(lr){3-4}\cmidrule(lr){5-6}
 & & Billow & \texttt{kite\_fem} & Billow & \texttt{kite\_fem} \\
\midrule
\texttt{La} & 0.848--0.903 & 2.4 & 2.4 & +1.9 & -0.6 \\
\texttt{Lb} & 0.979--1.280 & 16.5 & 10.6 & +16.5 & +8.3 \\
\texttt{Lc} & 0.903--1.276 & 18.8 & 14.3 & +18.8 & +9.7 \\
\texttt{Ld} & 0.749--1.259 & 24.5 & 20.8 & +24.5 & +20.8 \\
\texttt{Le} & 0.920--1.217 & 13.0 & 16.2 & +13.0 & +16.2 \\
\texttt{Lf} & 0.749--1.259 & 24.1 & 20.8 & +24.1 & +20.8 \\
\texttt{Lg} & 0.903--1.276 & 18.4 & 14.3 & +18.4 & +9.7 \\
\texttt{Lh} & 0.979--1.280 & 16.9 & 10.6 & +16.9 & +8.3 \\
\texttt{Li} & 0.848--0.903 & 2.4 & 2.4 & +1.9 & -0.6 \\
\texttt{b} & 4.867--8.271 & 30.3 & 22.9 & +30.3 & +22.9 \\
\texttt{LcsTL} & 4.029--4.888 & 10.3 & 6.8 & +10.3 & +5.1 \\
\texttt{LcsTR} & 4.029--4.888 & 10.4 & 6.8 & +10.4 & +5.1 \\
\bottomrule
\end{tabular}
\end{table}

\begin{figure}[htbp]
\centering
\includegraphics[width=\textwidth]{lengths.png}
\caption{The twelve measured quantities per load case, model over measurement.}
\end{figure}

\subsection{The solved kite}

\begin{figure}[htbp]
\centering
\includegraphics[width=0.94\textwidth]{canopy_1.png}\\[-1mm]
\includegraphics[width=0.94\textwidth]{canopy_5.png}
\caption{The solved shape under self-weight (top) and with the tips spread by
\SI{5}{kg} (bottom), with the measured markers overlaid. The canopy is shaded
by facet orientation, so the billow between struts and the sag of each bay are
visible directly. Both are the $EI\times0.36$ solution.}
\end{figure}

\begin{figure}[htbp]
\centering
\includegraphics[width=0.94\textwidth]{regime_1.png}
\caption{The same self-weight case with each canopy triangle coloured by its
tension-field state. This is a prediction the membrane formulation makes and a
tension-only spring net cannot: under self-weight almost the whole canopy is
wrinkled --- one principal stress relaxed to zero --- with taut material only
where the bridle pulls through, and a handful of fully slack triangles at the
tip.}
\end{figure}

\clearpage
\subsection{Geometry, case by case}

The figures that follow show the inflatable structure --- leading edge and the
ten struts --- in front, top and side projection, with the as-built shape for
reference. The canopy is omitted for legibility; it carries little of the load
in this case (Section~\ref{sec:what}).

Black dots are Haanen's photogrammetry markers, carried into the model frame by
a rigid transform. For these figures the registration uses the \emph{leading
edge alone}, so the two leading edges lie on top of one another and the
remaining disagreement appears where it physically is --- in the struts and the
canopy --- rather than being spread evenly over the structure by a global fit.
The error table in Section~\ref{sec:shape} uses its own per-model alignment over
all curves, which is the fair one for scoring; these overlays are for reading,
not for measuring. The markers stop short of the tips because the outermost bays
are not instrumented.

\begin{figure}[htbp]
\centering
\includegraphics[width=\textwidth]{geometry_1.png}
\caption{Load case 1.}
\end{figure}

\begin{figure}[htbp]
\centering
\includegraphics[width=\textwidth]{geometry_2.png}
\caption{Load case 2.}
\end{figure}

\begin{figure}[htbp]
\centering
\includegraphics[width=\textwidth]{geometry_3.png}
\caption{Load case 3.}
\end{figure}
\clearpage

\begin{figure}[htbp]
\centering
\includegraphics[width=\textwidth]{geometry_4.png}
\caption{Load case 4.}
\end{figure}

\begin{figure}[htbp]
\centering
\includegraphics[width=\textwidth]{geometry_5.png}
\caption{Load case 5.}
\end{figure}

\begin{figure}[htbp]
\centering
\includegraphics[width=\textwidth]{geometry_6.png}
\caption{Load case 6.}
\end{figure}
\clearpage

\begin{figure}[htbp]
\centering
\includegraphics[width=\textwidth]{geometry_7.png}
\caption{Load case 7.}
\end{figure}

\begin{figure}[htbp]
\centering
\includegraphics[width=\textwidth]{geometry_8.png}
\caption{Load case 8.}
\end{figure}

\begin{figure}[htbp]
\centering
\includegraphics[width=\textwidth]{geometry_9.png}
\caption{Load case 9.}
\end{figure}
\clearpage

\begin{figure}[htbp]
\centering
\includegraphics[width=\textwidth]{geometry_10.png}
\caption{Load case 10.}
\end{figure}


\subsection{The centre-load cases, which neither code reproduces}

Cases 2, 3, 7 and 8 put \SI{9.7}{} or \SI{25.2}{kg} on the mid-span leading
edge --- a node with \emph{no bridle attached}, carrying up to twice the
kite's own weight. Measurement has the span slightly \emph{widening} against
bare gravity (\SI{5.04}{} against \SI{4.87}{m}): the kite tilts forward and
the tips fold out.

Neither code reproduces that. \texttt{kite\_fem} overstates the span by
34--48\%, and its author is explicit about the cause: ``by applying such a
local load directly on the inflatable beam, the cross-section deforms and the
stiffness of the tube reduces \dots\ the simplification made by modelling such
a complex structure as a beam element cannot capture this'' (thesis \S7.1).
Billow fails differently and more visibly: it converges --- to residuals of
order \SI{1e-8}{N} --- onto a \emph{folded} branch, closing the span to under
\SI{1}{m}. Seeding from the gravity equilibrium and ramping the load in does
not recover the physical branch. The energy is non-convex, and this is a
second local minimum, not a numerical failure: a converged wrong answer.

These four cases are marked $\dagger$ in Table~1 and reported separately
rather than averaged in, because a mean over them measures how each code
fails rather than how well it works.

\clearpage
\section{What this case actually tests}
\label{sec:what}

The comparison was built to isolate the canopy formulation, since that is the
one element type that differs between the two codes. It does not do that. A
sensitivity sweep on case~1:

\begin{center}
\begin{tabular}{lr}
\toprule
change & span (m) \\
\midrule
baseline & 7.19 \\
canopy $Et$ \SI{5000}{}$\rightarrow$\SI{1000}{N/m} & 7.15 \\
canopy deleted entirely ($Et=\SI{1}{N/m}$) & 7.36 \\
tube moment--curvature curve scaled by $0.3$ & \textbf{3.73} \\
\bottomrule
\end{tabular}
\end{center}

Deleting the canopy outright moves the answer less than a quarter as far as a
modest change in tube bending. \textbf{This is a tube-bending test}, and any
canopy conclusion read off it would be unsupported.

\paragraph{Collapse.} None of the three codes models post-collapse behaviour.
Breukels' source data characterises it --- the measured moment falls past
collapse --- but \texttt{kite\_fem} computes a collapse flag and never acts on
it (a literal \texttt{TODO: verify collapse} in \texttt{BeamElement.py}), and
Billow saturates at $M_\text{max}$ instead of dropping. A law that does follow
collapse, fixed uniquely by requiring it to match both the fitted initial
stiffness and the fitted collapse curvature,
\begin{equation}
W(\kappa) = EI_0\,\kappa_c^{2}
\left[1-\left(1+\tfrac{\kappa}{\kappa_c}\right)e^{-\kappa/\kappa_c}\right],
\qquad
M(\kappa) = EI_0\,\kappa\,e^{1-\kappa/\kappa_c},
\end{equation}
\emph{converges without difficulty} (56--77 iterations, $\sim$\SI{3}{s},
residual \num{3e-9}) and changes the answer by \SI{20}{mm}. It changes nothing
here because Billow's beams reach only 31\% of the collapse curvature and none
of the 98 is flagged collapsed. Note also that collapse arrives at
$0.7\,\kappa_0$, i.e.\ at only $\sim$52\% of $M_\text{max}$, so the saturating
plateau is already too generous well before the cliff. It matters in general;
it does not matter in this dataset.

\section{Cost, convergence, and what the \texttt{kite\_fem} numbers are}

\begin{table}[htbp]
\centering
\caption{Solver cost. $^{\ast}$ marks cases that never reached
\texttt{kite\_fem}'s own \num{0.01} tolerance. \textbf{The two time columns were
measured on different machines} and are not a controlled benchmark --- see the
caveats below.}
\begin{tabular}{r rr rrr}
\toprule
LC & \multicolumn{2}{c}{\texttt{kite\_fem}} & \multicolumn{3}{c}{Billow} \\
\cmidrule(lr){2-3}\cmidrule(lr){4-6}
 & $t$ (s) & final tol. & $t$ (s) & iterations & $\|r\|$ (N) \\
\midrule
1 & 546 & 0.010 & 7 & 89 & 2e-07 \\
2 & 671 & 0.010 & 15 & 200 & 9e-08 \\
3 & 465 & 0.229$^{\ast}$ & 19 & 228 & 4e-07 \\
4 & 443 & 0.010 & 16 & 212 & 1e-07 \\
5 & 619 & 0.012$^{\ast}$ & 18 & 240 & 4e-07 \\
6 & 576 & 0.010 & 6 & 86 & 4e-09 \\
7 & 621 & 0.010 & 15 & 209 & 4e-09 \\
8 & 468 & 0.076$^{\ast}$ & 17 & 235 & 6e-09 \\
9 & 483 & 0.010 & 15 & 187 & 1e-07 \\
10 & 626 & 0.011 & 16 & 227 & 4e-07 \\
\midrule
mean & 552 & & 14 & 191 & \\
\bottomrule
\end{tabular}
\end{table}

Billow reaches force balance to better than \SI{5e-7}{N} on every case, against
a target of \SI{1e-4}{N}, in a mean of \SI{14}{s}. \texttt{kite\_fem}'s recorded
times average \SI{552}{s} and three of its ten cases stop above its own
tolerance (\num{0.229}, \num{0.076}, and a marginal \num{0.0117}).

Robustness required one addition to the solver. On slack-dominated
configurations a near-zero-curvature search direction is answered with a step of
order \num{1e4}, the objective overflows, and the line search collapses before
restoration can help. A move limit --- a trust region expressed as bounds IPOPT
already enforces --- cures this and reproduces the direct solve to the digit.
One caveat came with it: once bounds are active, IPOPT reports success for the
\emph{boxed} problem while the iterate sits on the boundary, so the acceptance
test must key off the force residual and never off the solver's own verdict.

\subsection{Caveats on every \texttt{kite\_fem} number in this report}
\label{sec:caveats}

These matter enough to state plainly, because the comparison is the weakest part
of this work and the reader should discount it accordingly.

\paragraph{It was never reproduced.} Every \texttt{kite\_fem} figure here comes
from its \emph{committed} \texttt{model\_results.csv} and saved \texttt{.npz}
states, not from a run performed for this study. One attempt was made; it
produced no output in over ninety minutes and was abandoned. So nothing here
distinguishes ``\texttt{kite\_fem} as published'' from ``\texttt{kite\_fem} as it
runs today''.

What \emph{was} verified is narrower but load-bearing: our length extractor
reproduces its published table from its own saved states to \SI{1e-9}{m} on all
ten cases (a committed test). The rescoring is therefore sound even though the
solve is not reproduced.

\paragraph{Its solver settings are tuning choices we did not replicate.} The
published run uses an \texttt{I\_stiffness} of 15, a tolerance of \num{0.01}, a
step limit, a relaxation schedule, and an outer loop that calls
\texttt{adapt\_stiffnesses} --- progressively stiffening any element above 1\%
strain, up to four times. None of that was re-derived or assessed here. A
different setting could move its answers by an unknown amount.

\paragraph{It is not solving the model it states.} Its converged beams sit at
about $0.23\times$ its own constitutive law, with negative bending stiffness in
three of the four cases inspected, and its anchor node belongs to no element so
its equilibrium is held by \texttt{I\_stiffness} rather than by a load path ---
its own solutions drift \SIrange{0.29}{0.34}{m} bodily. Comparing against it is
therefore comparing against a numerical state, not a clean reference solution.

\paragraph{``\texttt{kite\_fem}'' is not one thing.} The results used here are
from \texttt{main} at \texttt{630c8c0}. AWETrim's own environment installs the
branch \texttt{mass\_proportional\_regularisation} at \texttt{6edff17} --- named
for the very regularisation implicated above. Any statement in this report about
\texttt{kite\_fem} refers to the former.

\paragraph{The bridle is not matched in the headline comparison.}
\texttt{kite\_fem} ran the bridle exactly as stored. Our primary results apply
the Table~B.1 corrections. Scored on the stored bridle for a like-for-like
comparison, Billow gives 15.1\% against \texttt{kite\_fem}'s 12.4\%; on the
corrected bridle it gives 15.7\%. The difference the bridle makes overall is
therefore \SI{0.6}{} percentage points --- small, but the comparison is not
bridle-matched and should not be read as though it were.

\paragraph{The metric is ours, not theirs.} Its published
\texttt{shapecorrelation} is not used, for the reason given in Table~1's
caption. Numbers here will therefore not match those reported upstream.

\paragraph{The timings are not a benchmark.} Its recorded times are from an
HP~Zbook Power~G7; ours are from a different machine, running Python against a
compiled CasADi/IPOPT stack while \texttt{kite\_fem} runs Cython
\texttt{pyfe3d} elements. Read the ratio as an order of magnitude at most, and
note that the two codes are not even solving to the same tolerance --- a factor
of \num{1e3} separates the residuals.

\section{What stiffness the measurement implies}
\label{sec:lambda}

Rather than fit bridle lengths, the whole tube moment--curvature curve was
scaled by a single factor $\lambda$ --- ``\texttt{moment\_max}'',
``\texttt{bending\_stiffness}'' and ``\texttt{torsion\_c1}'' together, so
$\kappa_0 = M_\text{max}/EI_0$ is unchanged and the curve keeps its shape while
only its magnitude moves. The two bare-gravity cases, the only two that
genuinely load the tubes, then say what stiffness the measurement wants:

\begin{center}
\begin{tabular}{rrrrr}
\toprule
$\lambda$ & \multicolumn{2}{c}{case 1, \SI{0.15}{bar}} & \multicolumn{2}{c}{case 6, \SI{0.25}{bar}} \\
\cmidrule(lr){2-3}\cmidrule(lr){4-5}
 & span (m) & error (\%) & span (m) & error (\%) \\
\midrule
1.00 & 7.276 & 25.9 & 7.468 & 19.5 \\
0.55 & 6.156 & 14.6 & 6.589 & 11.5 \\
0.45 & 5.569 &  9.3 & 6.120 &  7.7 \\
\textbf{0.35} & \textbf{4.674} & \textbf{8.3} & \textbf{5.376} & \textbf{5.5} \\
0.25 & 3.376 & 18.6 & 4.141 & 14.5 \\
\midrule
\multicolumn{1}{r}{measured} & 4.867 & & 5.453 & \\
\bottomrule
\end{tabular}
\end{center}

\noindent
Both cases imply $\lambda \approx 0.36$, and --- this is the part that matters
--- they imply the \emph{same} factor at both pressures. A pressure-dependent
error in the fit would not do that; a constant offset would. It is also further
evidence against the low-pressure extrapolation of $M_\text{max}$ discussed
above, which would have required more softening at \SI{0.15}{bar} than at
\SI{0.25}{bar}.

Run across all ten cases, a single $\lambda = 0.36$ gives a mean relative error
of 8.0\% against \texttt{kite\_fem}'s 12.4\%, and 6.2\% against 11.5\% on the two
cases that test the tubes. It also repairs the centre-load cases, which the
unscaled model got worst: case~2 falls from 27.6\% to 9.6\%, case~7 from 19.7\%
to 7.5\%. The tip-load cases barely move (6.5\%$\to$6.6\%, 5.2\%$\to$5.2\%),
exactly as the previous section predicts --- they are bridle-determined and
carry no information about the tube law.

\paragraph{This is an identification, not a fit.} One scalar, applied uniformly
to all 98 tube elements, calibrated on two cases and improving eight. It is
reported because a consistent, pressure-independent factor is evidence about the
fit rather than a free parameter absorbing error --- but it is not adopted as a
model change. Two things must be understood first.

\emph{Where does $0.36$ come from?} Inverting the fitted $EI_0$ through
$EI = E t \pi r^3$ gives an implied $E t$ of \SIrange{133}{155}{kN/m} for the
leading-edge tubes, nearly constant across $d = \SIrange{0.113}{0.202}{m}$ ---
so the fit's radius dependence is physically consistent, and the implied fabric
(\SIrange{0.5}{1.2}{GPa} at \SIrange{0.15}{0.25}{mm}) is plausible to soft for
polyester. The fit is therefore not obviously twice too stiff on material
grounds, which makes a calibration or usage mismatch the more likely
explanation.

Two candidates are worth checking against Breukels' thesis directly. First, his
tubes are discretised as \emph{rigid cylindrical segments joined by torque
springs}, and the $C_1$--$C_8$ coefficients were fitted so that such a chain
reproduces the measured beam --- a discretisation-dependent calibration, not an
intrinsic moment--curvature law, and re-using it as a continuum law could carry
a systematic factor. Second, his nomenclature defines an \emph{effective}
internal pressure $p_\text{eff}$ distinct from the gauge pressure; if the fits
are calibrated on $p_\text{eff}$, feeding \SI{0.15}{bar} gauge straight in is
wrong at exactly this kind of magnitude.

\section{How weakly the published data constrains a model}

The twelve published quantities are less information than they look. Every
left/right pair is identical to $0.0\times10^{0}$ --- \texttt{La}$\equiv$%
\texttt{Li}, \texttt{LcsTL}$\equiv$\texttt{LcsTR} --- so the data is exactly
mirrored and only six numbers are independent. The thesis itself scores just
three: the trailing-edge length (the sum of \texttt{La}--\texttt{Li}), the span,
and the tip-to-centre distance.

More limiting is what is absent. There is no height, no sag depth, and no
chordwise measure of any kind: every published quantity lies in the spanwise /
trailing-edge geometry. A model can match all three scored lengths and still
have the wrong shape --- the wrong camber, the wrong chordwise twist, the wrong
depth of sag --- with nothing in the dataset able to detect it. That is a real
ceiling on what any validation against this file can claim, including this one.

The information exists. The measurement was stereoscopic, giving full
three-dimensional marker positions, and the thesis plots measured geometry
against model geometry in top, side and front projection for all ten load cases
(its Figure~7.2). Only the twelve scalars were published.

\paragraph{The single most valuable thing to obtain} is therefore the 3D marker
coordinates behind Figure~7.2. They would turn this from a comparison of three
effective scalars into a comparison of shapes, and would immediately settle
questions the present data cannot touch --- among them whether Billow's wing is
uniformly too open or wrong in a particular region, and whether the centre-load
cases fail in the way the folded branch suggests.

\section{The proper test: 3D shape against the markers}
\label{sec:shape}

Everything above is scored on twelve numbers that reduce to six independent
ones. Haanen's marker clouds --- 62 to 76 three-dimensional points per load case
on the leading edge, the eight struts and the canopy --- allow the question that
data cannot answer: is the \emph{shape} right?

\paragraph{Method.} Markers are in the stereo rig's frame, so a \emph{rigid}
transform is needed --- rotation and translation only, never scale, since scale
is what is being tested. Correspondence comes from structure rather than
fitting: the measured struts are ordered across the span with their first marker
at the leading edge, so their roots pair with the model's eight strut
leading-edge nodes, giving eight anchors for a Kabsch fit. Both spanwise
orderings are tried, because the kite is near-symmetric and the marker numbering
direction is undocumented. The fit is then refined by ICP against the model
curves, so no marker is assumed to sit on a node. Reported errors are
point-to-\emph{surface}: the markers are stuck to the outside of the tubes while
the model polyline is the tube axis, and ignoring that radius --- 55 to
\SI{100}{mm} here --- would put a systematic offset into every case.

\begin{table}[htbp]
\centering
\caption{Point-to-surface RMS between the solved shape and the photogrammetry
markers. ``As built'' is the undeformed geometry, included as the null model:
anything not beating it has explained nothing.}
\begin{tabular}{rl rrr}
\toprule
LC & load & \multicolumn{2}{c}{Billow (mm)} & as built (mm) \\
\cmidrule(lr){3-4}
 & & EI x1.0 & EI x0.36 & \\
\midrule
1 & gravity & 480 & 100 & 695 \\
2 & centre 9.7 kg & 501 & 152 & 654 \\
3 & centre 25.2 kg & 507 & 171 & 639 \\
4 & tip 2 kg & 144 & 136 & 141 \\
5 & tip 5 kg & 101 & 92 & 102 \\
6 & gravity & 387 & 101 & 554 \\
7 & centre 9.7 kg & 399 & 151 & 525 \\
8 & centre 25.2 kg & 409 & 145 & 516 \\
9 & tip 2 kg & 119 & 111 & 114 \\
10 & tip 5 kg & 102 & 94 & 101 \\
\midrule
\multicolumn{2}{r}{mean, gravity only} & \textbf{434} & \textbf{101} & 625 \\
\multicolumn{2}{r}{mean, centre load} & \textbf{454} & \textbf{155} & 583 \\
\multicolumn{2}{r}{mean, tip load} & \textbf{117} & \textbf{108} & 115 \\
\multicolumn{2}{r}{mean, all ten} & \textbf{315} & \textbf{125} & 404 \\
\bottomrule
\end{tabular}
\end{table}

\paragraph{The scaled model reaches the measurement floor.} On the two cases
that load the tubes, the error falls from \SI{480}{} to \SI{100}{mm} and from
\SI{387}{} to \SI{101}{mm} --- onto the same $\sim$\SI{100}{mm} the tip cases
sit at, which is the marker placement and alignment floor. Across all ten the
mean falls from \SI{315}{} to \SI{125}{mm}, against \SI{404}{mm} for doing
nothing. That the factor identified from two scalar spans also fixes the full
three-dimensional shape, including the centre-load cases it was not fitted on,
is the strongest evidence in this report that it is a real property of the tube
law rather than a number absorbing error.

\paragraph{The tip-load cases are confirmed uninformative.} Their errors are
\SI{141}{}, \SI{102}{}, \SI{114}{}, \SI{101}{mm} for the \emph{undeformed}
geometry against \SI{144}{}, \SI{101}{}, \SI{119}{}, \SI{102}{mm} for the
unscaled model and \SI{136}{}, \SI{92}{}, \SI{111}{}, \SI{94}{mm} for the scaled
one --- all within about \SI{10}{mm} of each other. Doing nothing scores as well
as either model. This is the same conclusion the bridle-tension argument reached
in Section~8, now established directly in three dimensions rather than inferred,
and it is why the six-case and ten-case means elsewhere in this report should be
read with care.

\section{Summary}

\begin{itemize}
\item The distributed reference model is ill posed --- its anchor belongs to no
      element --- and its published results are held up by a numerical
      regularisation. Any future use of this dataset has to supply a suspension.
\item \texttt{kite\_fem}'s beams converge to $\sim0.23\times$ its own
      constitutive law, with negative bending stiffness in several cases. Its
      closer agreement with the measurement is not evidence of better physics.
\item The dataset is a tube-bending test --- and only two of its ten cases even
      are. It cannot decide a canopy model, which was the original reason for
      building the comparison. The 3D markers confirm this directly: on the
      tip-load cases the undeformed geometry scores as well as either model.
\item Against the 3D markers, $\lambda = 0.36$ takes the mean point-to-surface
      error from \SI{315}{} to \SI{125}{mm} (null model \SI{404}{mm}), and
      brings the two tube-loaded cases onto the $\sim$\SI{100}{mm} measurement
      floor --- having been identified on two scalar spans, not on the shape.
\item Post-collapse behaviour is unmodelled in all three codes and does not
      matter here or in coupled use, because the tubes never leave their initial
      slope.
\item The residual disagreement lives in $EI_0$. A single pressure-independent
      factor $\lambda \approx 0.36$, identified on the two gravity cases,
      improves eight of the ten and beats \texttt{kite\_fem} overall
      (8.0\% against 12.4\%). It is reported as an identification, not
      adopted: two specific mismatches in how the Breukels fit is being
      re-used --- its torque-spring discretisation and its effective
      pressure --- should be checked against the source first.
\end{itemize}

\section*{Data and code}

Measurements, the reference model and \texttt{kite\_fem}'s published results are
vendored under \texttt{data/LEI-V3-KITE/hanging\_validation/} from
\texttt{awegroup/kite\_fem} (MIT, \textcopyright~2025 Patrick Roeleveld),
\texttt{main} at \texttt{630c8c0}. The model builder is
\texttt{scripts/structural/hanging\_kite.py}; the driver, plots and this
report's tables are the neighbouring \texttt{run\_}, \texttt{plot\_} and
\texttt{make\_hanging\_report.py} scripts.

Thesis references are to P.\,I.\,H.~Roeleveld, \emph{Fast finite element
modelling of bridled leading-edge inflatable kites}, MSc thesis, TU~Delft,
February 2026.

\end{document}
